% Chapter 10: Parallel and Delegated Workflows
\chapter{Parallel and Delegated Workflows}
\label{ch:parallel}

The natural mode of Claude Code is sequential: one conversation,
one context, one task. But many real workflows benefit from parallelism---
running multiple investigations simultaneously, reviewing code from
multiple perspectives, or working on independent modules concurrently.

\section{Parallel Session Strategies}
\label{sec:parallel-strategies}

\marginnote[Official]{Desktop app, web, and Agent Teams for simultaneous work. Git worktrees (\code{--worktree}) for isolated parallel sessions.\sourceurl{https://code.claude.com/docs/en/common-workflows}}

\subsection{Git Worktrees}

\official{The \code{--worktree} flag creates an isolated copy of your repository
for parallel work.}

\begin{minted}{bash}
# Start parallel work on two features
claude --worktree "Implement caching layer"
claude --worktree "Add rate limiting"
\end{minted}

Each worktree:
\begin{itemize}
  \item Gets its own git branch
  \item Has a complete, isolated copy of the repository
  \item Can be merged independently when complete
  \item Is cleaned up automatically if no changes are made
\end{itemize}

\subsection{Multi-Perspective Review}

\official{Writer/Reviewer pattern}: run one session to make changes,
a separate session to review them. Neither session's context contaminates the other.

\practitioner{Extend this to multi-perspective review with subagents:}

\begin{enumerate}
  \item \textbf{Correctness agent}: ``Does this code do what it claims?''
  \item \textbf{Security agent}: ``Are there vulnerabilities?''
  \item \textbf{Performance agent}: ``Are there efficiency concerns?''
  \item \textbf{Style agent}: ``Does this match project conventions?''
\end{enumerate}

Running four agents in parallel takes the same wall-clock time as one.

% -------------------------------------------------------------------
\section{When to Delegate vs.\ Do Directly}
\label{sec:delegation-decision}

\practitioner{Not everything benefits from delegation.
The overhead of spawning a subagent is justified only when the task
is independent and context-heavy enough to warrant isolation.}

\subsection{Delegate When}

\begin{itemize}
  \item The task is independent of your current work
  \item It requires heavy context that would pollute your main session
  \item It can run in the background while you continue other work
  \item You need an unbiased second opinion
  \item Multiple tasks can run in parallel for wall-clock savings
\end{itemize}

\subsection{Do Directly When}

\begin{itemize}
  \item The task is quick (under 2 minutes)
  \item It depends on your current session's context
  \item Sequential reasoning is required (step B depends on step A's result)
  \item The task is too simple to justify the overhead
\end{itemize}

\subsection{Parallel Validation Example}

\practitioner{A common high-value pattern: parallel pre-commit validation.}

Instead of running lint, tests, type checking, and coverage sequentially
(20 minutes), delegate each to a subagent and run in parallel (5 minutes):

\begin{minted}{text}
Parallel agents:
  Agent 1: Run linter, report violations
  Agent 2: Run test suite, report failures
  Agent 3: Run type checker, report errors
  Agent 4: Run coverage analysis, report gaps
\end{minted}

All four results return to your main session.
You address the issues in priority order.

\begin{keyinsight}
Parallelism in Claude Code is not about speed alone---it is about
\emph{perspective isolation}. A reviewer who cannot see the writer's
reasoning catches different bugs. A security agent that only sees code
(not intent) finds different vulnerabilities. The diversity of perspectives
is the value, not just the time savings.
\end{keyinsight}
